\section{代数幾何と学習理論からの抜粋}

サンプルから真の構造を推測する学習モデルにおいては，真の構造はパラメータ空間の中の特異点を持つ解析的集合になっている．

（意味がわからん．真の構造とは，真の分布の構造のことなのか？) ここでは，その場合に必要となる学習理論の基本定理を紹介する．


\subsection{統計的推測の問題設定}
ユークリッド空間 \({\mathbf{R^N}}\) 上の確率密度関数 \({q(x)}\) があって，確率分布 \({q(x) dx}\) に従う\({n}\) 個の独立な確率変数
\(\mathbf{X^N}=(X_1, X_2,..., X_N)\) があるとする．
このとき \({\mathbf{X^N}}\) をサンプルという．（学習データ，例とか）
これの生成源となる確率分布 \({q(x) dx}\) を真の分布とか情報源という．

学習理論（統計的推測）でやりたいことは，
「サンプル\({\mathbf{X^N}}\)
から真の分布 \({q(x)}\) を推測する」
である．


ユークリッド空間上の確率密度関数 \({q(x)}\)（ってなに？）には無限の可能性があるので，
一切の仮定をおかずに，有限個のサンプルを見ているだけでは推測はできない．
現実的に推測を可能にするためには何らかの家庭を置く必要がある．
ここで置く仮定によって以下で説明するベイズ学習（ベイズ推論）が可能になる．


\begin{quote}
なので，ベイズ法とは，高次の意味での学習すなわち真の分布の推測を達成するための一連の仮定の置き方に関するルールと，それらの仮定によって可能になる方法論のあつまり（フレームワーク）と捉えられる？
\end{quote}



\subsection{ベイズ法}
データ分析者は，パラメータ \({w \in \mathbf{W^d}}\) によって定まる\({x \in \mathbf{R^N}}\) の条件付き確率密度分布 \({p(x \mid w) }\) と，パラメータ \({w \in \mathbf{W^d} }\) の確率密度分布  \({\phi(w)}\) を定める．
\({p(x\mid w)}\) を統計モデルや確率モデル，\({\phi(w)}\) を事前分布という．

次に，事後分布を以下で定義する．
